\documentclass[11pt,a4paper]{article}

\usepackage[utf8]{inputenc}
\usepackage[T1]{fontenc}
\usepackage[spanish]{babel}
\usepackage[margin=2.5cm]{geometry}
\usepackage{csquotes}
\usepackage{graphicx}
\usepackage{booktabs}
\usepackage{amsmath}
\usepackage{listings}
\usepackage[hidelinks]{hyperref}
\usepackage[backend=biber,style=ieee]{biblatex}
\addbibresource{referencias.bib}

\begin{document}

\begin{titlepage}
  \centering
  {\large Universidad Técnica Estatal de Quevedo\par}
  {\large Facultad de Ciencias de la Computación\par}
  {\large Carrera de Software\par}
  \vspace{0.5cm}
  {\large Aplicaciones Distribuidas (ISR-701) \,--\, 7mo nivel, paralelo A\par}
  \vspace{2cm}
  {\LARGE\bfseries Práctica experimental de la Unidad 1:\\
  comunicación entre procesos distribuidos\\con sockets TCP y gRPC\par}
  \vspace{1.5cm}
  {\large Código de la actividad: GA-SUM-02 / PE-U1\par}
  \vspace{1.5cm}
  {\large\bfseries Grupo GLMR\par}
  \vspace{0.4cm}
  \begin{tabular}{c}
    García Bazurto Kevin Leonel \\
    Lombeida Escaleras Bryan Humberto \\
    Manotoa Álava Carlos Alexander \\
    Reinoso Vélez Eduardo David
  \end{tabular}
  \vspace{1.5cm}

  {\large Docente: Ing. Gleiston Guerrero Ulloa, PhD\par}
  \vspace{0.4cm}
  {\large Noviembre de 2026\par}
  \vspace{1.5cm}
  {\large Repositorio:\par}
  {\small\url{https://github.com/yaser-uwu/GA-SUM-02-PE-U1}\par}
\end{titlepage}

\section{Objetivo}

Implementar y comparar dos mecanismos de comunicación entre procesos, sockets TCP y
gRPC con Protocol Buffers, en un sistema de tres nodos que corre en \texttt{localhost}.
Se mide la latencia de ida y vuelta y el rendimiento con 100 envíos por corrida, y se
comprueba con relojes lógicos de Lamport el orden de los eventos entre los nodos.

\section{Hipótesis}

Se espera que los sockets TCP tengan menor latencia que gRPC. En gRPC cada mensaje pasa
por trabajo adicional: se codifica en binario con Protocol Buffers y viaja sobre una capa
más de protocolo, HTTP/2. En TCP solo se envía el JSON.

Esta hipótesis se escribió el 18 de noviembre de 2026, antes de ejecutar las mediciones
oficiales.

\end{document}
